%!TEX TS-program = lualatex
%!TEX encoding = UTF-8 Unicode

\documentclass[12pt, addpoints]{exam}

%\printanswers


\usepackage[left=1in,right=0.75in,top=1in]{geometry}     

\usepackage{fontspec}
\def\mainfont{Linux Libertine O}
\setmainfont[Ligatures={TeX, Common}, BoldFont={* Bold}, ItalicFont={* Italic},  Numbers={Proportional, OldStyle}]{\mainfont}
%\setmonofont[Scale=MatchLowercase]{Inconsolata} 
\setsansfont[Scale=MatchLowercase]{Linux Biolinum O} 
\usepackage{microtype}


% This defines \amper for the fancy ampersand
% to be used in the header. See
% https://tex.stackexchange.com/a/58185/39194
\usepackage{xspace}
\newfontfamily\amperfont[Style=Alternate]{Linux Libertine O}
\makeatletter
\DeclareRobustCommand{\amper}{{\amperfont\ifx\f@shape\scname\smaller[1.2]\fi\&}\xspace}
\makeatother

\usepackage[parfill]{parskip}

\usepackage{graphicx}
	\graphicspath{%
	{/Users/goby/Pictures/teach/434/exams/}}%

% Used to get the last two digits of the year for the header below.
\def\short#1{\csname @gobbletwo\expandafter\endcsname\number#1}

\pagestyle{headandfoot}
\firstpageheader{\textsc{bi}\,434/634 Marine Ecology \amper Conservation, Exam 3, \textsc{f}\short{\year}}{}{\numpoints\ points}
\runningheader{}{}{\small{pg. \thepage}}
%\runningheadrule

\footer{}{}{}%{\makebox[0.75in]{\hrulefill}}
\extrafootheight{-0.5in}

\usepackage{enumitem}
\setlist{leftmargin=*}
\setlist[1]{labelindent=\parindent}

\usepackage{hanging}

\usepackage[sc]{titlesec}

\pointsinmargin
\marginpointname{ pts}

\usepackage{multicol}

%% Remove comment %% to print answer key.
\unframedsolutions
\renewcommand{\solutiontitle}{}
\SolutionEmphasis{\bfseries}

\renewcommand{\questionshook}{%
	\setlength{\leftmargin}{-\leftskip}%
}


\newcommand*\AnswerBox[2]{%
    \parbox[t][#1]{0.92\textwidth}{%
    \begin{solution}#2\end{solution}}
    \vspace{\stretch{1}}
}

\newenvironment{AnswerPage}[1]
    {\begin{minipage}[t][#1]{0.92\textwidth}%
    \begin{solution}}
    {\end{solution}\end{minipage}
    \vspace{\stretch{1}}}


\newcommand{\fst}{$F_{ST}$}
%\newlength{\basespace}
%\setlength{\basespace}{5\baselineskip}



\begin{document}

Each question is based on a chapter from the text and one or more recent scientific publications. 
The publications are available from the course Moodle page. Read each question, find the 
relevant material in the chapter and publication, and then answer the question as fully
as possible.

Most questions have a \emph{technically correct} answer but I am more interested in
the depth and justification of your answers. I expect that you will explain your reasoning 
for your answers. You may quote directly (use quotation marks to  show when you are quoting) from the text and papers to 
support your reasoning but do not quote extensively. The majority of writing must be in
your own words. \textbf{When you quote or otherwise reference the text, provide specific page, figure, or table numbers.}

Upload your answers to the drop box by the due date and time.

\begin{questions}

\begin{EnvFullwidth}
	\subsubsection*{Chapter 12: Ecology of seagrass communities}
	
	In Chapter 12, Duffy, Hughes, and Moksnes (2014) showed how small
	mesograzers can regulate ecosystem function in seagrass beds, through
	the mutualistic mesograzer model (read that three times fast!). 
\end{EnvFullwidth}

\question
Facilitation cascades may also improve seagrass ecosystem function. Zhang and Silliman (2019) tested this hypothesis with \textit{Zostera} (a seagrass genus) and \textit{Atrina rigida}, the pen clam. The authors \emph{removed} mesograzers to determine whether the pen clams alone improved ecosystem function.


\begin{parts}
	
	\part[5]
	In your own words, describe what is a facilitation cascade? Again, in your words, what is the facilitation cascade proposed by this study?
	
	\ifprintanswers \textbf{%
		Faciliation cascades are a series of positive interactions that promote
		diversity. Here, they test whether \textit{Zostera} facilitates the
		settlement and survival of a second foundation species, \textit{Atrina.}
	}\fi
	
	
	\part[6]
	Use Figs.\ 2–5 in any combination you deem necessary to answer these brief questions: i) Does \textit{Zostera} seem to facilitate the presence of \textit{Actina?} ii) Does the presence of \textit{Actina} improve overall biological diversity in the ecosystem? iii) Does the presence of \textit{Actina} seem to improve ecosystem function?
	
	\ifprintanswers \textbf{%
		i) Fig.\ 2 shows that pen clam density was higher and mortality lower in seagrass beds compared to surrounding sand flats. ii) Fig.\ 5 showed that abundance and SW diversity increased, but iii) ecosystem function did not improve, at least by the measured used.
	}\fi
	
	\part[4]
	Do you think that long-term facilitation between \textit{Zostera} and
	\textit{Actina} could improve seagrass ecosystem function? Why or why not? Provide a well-reasoned answer.
	
	\ifprintanswers \textbf{%
		Open-ended but I kinda think function could be improved. The clams may
		help to stabilize the sediments for the seagrass. Pseudofaeces may 
		further provide nutrients. Filter feeding may contribute to clearing. 
		The hard surfaces may facilitate growth of other primary producers.
	}\fi
	
	
\end{parts}

\begin{EnvFullwidth}
	\subsubsection*{Chapter 11: Salt marsh communities}
	
	In chapter 11, Bertness and Silliman (2014) discussed how top-down 
	procceses influence diversity and ecosystem function in salt marsh 
	communities. Further, in lecture, I have stressed how local disturbances
	can increase diversity over time at regional scales.
	
\end{EnvFullwidth}

\question
Elschot et al.\ (2017) studied whether the local foraging affects of
Greylag Goose affected top-down vs bottom-up effects in salt marshes
at different scales. They examined both geographic and temporal scales. Study the paper to learn the details of the scales they considered. (Notes: i) grubbing is digging below the surface to feed. ii) the arrow colors given in the legend for Fig.~2 are not correct. Replace white and yellow arrow color, respectively, with black and white arrows color.)

\begin{parts}
	
	\part[5]
	What was the specific top-down process studied? Why is this a top-down process? What was the specific bottom-up process studied? Why is this a bottom-up process?
	
	
	\ifprintanswers \textbf{%
		Top-down was grubbing by the geese, which opens bare batches in the vegetation. Here, the geese are controlling the plants, so this is top-down. Bottom-up is soil accretion, or the accumulation of soil above sea level. Soil is necessary for plant growth so this is bottom up.
	}\fi
	
	\part[6]
	What is the dominant plant species in these salt marshes, prior to 
	disturbance by the geese? Fig.~5 suggests three stages of succession as bare patches created by the geese return to the dominate vegetation state: early (2 years), middle (6 years), and late (12 years) as the marsh returns to the dominant species. Describe the 
	pattern of succession by listing the dominant pairs of species present 
	from Fig.~5 at each of the three stages. Use percent vegetation cover 
	from Table~1 to support your answer. You do not need to discuss other 
	species from Table 1 that are not included in Fig.~5.
	
	\ifprintanswers \textbf{%
		Early species are \textit{Pucinella} and \textit{Aster}. Their 
		abundances peak at year 2 at 11 and 28\% respectively. Middle species
		are Glaux and Salicornia at 6 and 13\%, respectively. NOTE to MST: other species have higher percent cover. Finally, \textit{Elytrigia} and \textit{Bolboschoenus} dominate the late stage with 15 and 38\%.
	}\fi
	
	\part[4]
	How does the presence of the geese increase diversity of the salt marsch at the scaled defined in this study?
	
	\ifprintanswers \textbf{%
		The bare patches allow new plant species to colonize. During succession, each patch will have different plant species, depending on when the patches were made. This increases diversity over time. Further, as patches are created across the marsh, the patches will be in different states of recovery, thus increasing diversity across the region.
	}\fi
	
\end{parts}


\subsubsection*{Chapter 13: Coral reef ecosystems}

\fullwidth{Côté and Knowlton (2014) discuss several stressors that 
can negatively affect the health state of coral reefs. They also discuss
briefly prospects for conservation driven by these stressors. One prospect
that they do not discuss is refugia from stressors.}

\question
Glynn~(1996) discussed the potential effects of global climate change
on shallow-water coral reefs. 

\begin{parts}
\part[10]
Glynn~(1996) suggested that deep (low-light, or “mesophotic”) 
reefs might serve as refuges for some shallow-water coral 
species. Explain, with justification, the \emph{specific} stressors 
listed by Côté and Knowlton (2014) that you think would be 
ameliorated for shallow species if they were able to survive 
in mesophotic reefs. 

\ifprintanswers \textbf{%
	Specific stressors mentioned by Côte and Knowlton include overfishing and destructive fishing, poor water quality, invasive species, warming, and acidification. Warming is the best one to discuss and must be included. Accept any reasonable argument for the second one. 
}\fi


\part[5]
Glynn~(1996) suggested that the value of mesophotic refuges 
might vary among species with different growth forms and 
colony sizes. What are the two \emph{specific} forms that Glynn 
mentions. Give a reasonable justification for his choices; 
that is, why would one form be more likely to survive in 
a deep refuge than the other form? Be clear and specific.

\ifprintanswers \textbf{%
Glynn specifically mentions massive coral species of a large colony size with slow skeletal growth and branching coral species of a small to moderate colony size with a more rapid growth. He cites a previous study that showed branched corals bleached more readily than massives. Massives were more likely to bleach on top than the sides, so they might maintain productivity.
}\fi

\end{parts}


\question
Bongaerts~et~al.~(2017) tested the deep-reef refuge 
hypothesis for two coral species, \textit{Agaricia fragilis}
and \textit{Stephanocoenia intersepta.}

\begin{parts}
\part[5] 
Are these two species a good test of Glynn's hypothesis
regarding coral growth forms? Why or why not? \textsc{Hint:} Google is your friend.

\ifprintanswers \textbf{%
\textit{Agaricia} has a plate growth form which Glynn suggests is more susceptible to bleaching than massive. \textit{Stephanocoenia} is similar to massive corals in shape but not massive in the sense of big.
}\fi

\part[5]
Do the results from Bongaerts~et~al.\ (2017) for each coral
species support or falsify the deep-reef refuge hypothesis?
Use specific examples from the paper to support your conclusions.
\textsc{Hint:} Supplementary figure \textsc{s2} and table \textsc{s3} could be helpful.

\ifprintanswers \textbf{%
\textit{Agaricia} shows greater genetic differences between deep and shallow pops than within suggesting that deep to shallow gene flow cannot be relied upon. In contrast, \textit{Stephanocoenia} showed complete admixture, suggesting that deep populations could serve as a refuge source.} \fi

\end{parts}

\subsubsection*{Chapter 17: deep-sea hydrothermal vent communities}

Mullineaux (2014) addressed concerns about mining of hydrothermal vent communities but little data on actual effects are available  so we'll take a different approach for this chapter. 

Lique and Thomas (2018) used complex models to show how the Atlantic component of the Ocean Conveyor (the \textsc{amoc}) will be affected by large increases in atmospheric CO\textsubscript{2}. Their results are \emph{dense} but summarized in a perspective piece by Tamsitt (2018).  Read Tamsitt and compare her Fig.~1 to Fig.~5 of Lique and Thomas (do \emph{not} read all of Lique and Thomas unless you want to; just compare the figures.) Their results suggest the \textsc{amoc} will weaken, meaning the rate that water mass moves will slow-down. Weakening also might mean that water mass will not sink as deep below the surface.

\question
In a separate study, Turner and colleagues (2019) identified potential ecosystem functions and services from hydrothermal vents. See especially their Table~2 and Fig.~1. 

\begin{parts}
%\part[5]
%If the conveyor is weakening (slowing down, not going as deep), describe how you think this will affect hydrothermal vent communities. I am asking you to write in essence a hypothesis. Be sure it is robust and well-justified. You should address at least oxygen and nutrient availability.

\part[5]
Pick \emph{one} ecosystem function identified by Turner~et~al.~(2019) and describe how you think it might be affected by a weakening \textsc{amoc}. I am asking you in essence to write a hypothesis. Be sure it is robust and well-justified.

\ifprintanswers \textbf{%
Open-ended. Accept anything reasonable.} \fi

\part[5]
If the \textsc{amoc} is weakening, describe how the potential change of a hydrothermal regulating service (Fig.~1 of Turner et~al.) might affect shallow-water marine ecosystems. Think carefully about the potential roles described for the regulating service. As well, Tamsitt mentions a specific function of the \textsc{amoc} that fits under the broader umbrella listed by Turner et al.

\ifprintanswers \textbf{%
	Open-ended. Accept anything reasonable.} \fi

\end{parts}

\subsubsection*{Chapter 19: Climate change and marine communities}

Climate change is addressed, directly or indirectly, in most chapters of your text but Bruno, Harley, and Burrows (2014) discuss the many possible negative affects of climate change across marine ecosystems. Can there be positive effects of climate change on any marine ecosystem?

\question
Sunday and her colleagues (2017) tested for the possible effects of acidification on four biogenic habitat types. 

\begin{parts}
\part[10]
List the four biogenic habitats. For each habitat type, describe how acidification affects biodiversity in that habitat. Provide specific evidence from their results for each habitat.

\ifprintanswers \bfseries{%
1. Coral reef. Habitat complexity is reduced by shifting to corals with less complex morphology. Reduced complexity reduces available niches, lowering diversity.

\smallskip

2. Mussel beds. Large \textit{Mytilus} species will be replaced by smaller species, reducing structural complexity for supporting interstitial fauna, the organisms that grow between the mussels.

\smallskip

3. Seagrasses \newline
4. Fleshy erect macroalgae: most often, biomass increased for these two foundation spp, thus increasing overall diversity. Presumably, this is due to increased CO2. A few scenarios suggest other photosynthetic competitors increased while seagrass or macroalgae decreased.
} \fi

\part[5]
Describe what \emph{you} think are the implications for global biodiversity, using their results to guide your reasoning. You may include ecosystems other than the four studied above, but tell how they relate to one or more of the four studied habitat types.

\ifprintanswers \textbf{%
	Open-ended. Accept anything reasonable.} \fi


\end{parts}
\end{questions}

\subsubsection*{Literature cited}

\begin{hangparas}{\leftmargin}{1}


%Angelini, C., J.\,N.\,Griffin, J.\,van de Koppel, L.\,P.\,M\,Lamers, A.\,J.\,P.\,Smolders, M.\,Derksen-Hooijberg, T.\,van der Heide, and B.\,R.\,Silliman. 2016. A keystone mutualism underpins resilience of a coastal ecosystem to drought. Nature Communications~7: 12473 | DOI: 10.1038/ncomms12473


Bangaerts, P., C.\,Riginos, R.\,Brunner, N.\,Englebert, 
S.\,R.\,Smith, and O.\,Hoegh-Guldberg. Deep reefs are not 
universal refuges: Reseeding potential varies among coral species.
Science Advances 3: e1602373.

Elschot, K., A.\ Vermeulen, W.\ Vandenbruwaene, J.\ P.\ Bakker, T.\ J.\ Bouma, J.\ Stahl, H.\ Castelijns, and S.\ Temmerman. 2017. Top-down vs. bottom-up control on vegetation composition in a tidal marsh depends on scale. PLoS ONE 12(2): e0169960. doi:10.1371/journal.pone.0169960

Glynn, P.\,W. 1996. Coral reef bleaching: facts, hypotheses and 
implications. Global Change Biology 2: 495–509.

Lique, C. and M.\,D.\,Thomas. 2018. Latitudinal shift of the Atlantic Meridional Overturning Circulation source regions under a warming climate. Nature Climate Change 8: 1013–1020. DOI: 10.1038/s41558-018-0316-5

Sunday, J.\,M., K.\,E.\,Fabricius, K.\,J.\,Kroeker, K.\,M.\,Anderson, N.\,E.\,Brown, J.\,P.\,Barry, S.\,D.\,Connell, S.\,Dupont, B.\,Gaylord,
J.\,M.\,Hall-Spencer, T.\,Klinger, M\,Milazzo, P.\,L.\,Munday, B.\,D.\,Russell, E.\,Sanford, V.\,Thiyagarajan, M.\,L.\,H.\,Vaughan, S.\,Widdicombe, and C.\,D.\,G.\,Harley. 2017. Ocean acidification can mediate biodiversity
shifts by changing biogenic habitat. Nature Climate Change 7: 81–85. DOI: 10.1038/nclimate3161.

Tamsitt, V. 2018. Moving windows to the deep ocean. Nature Climate Change 8: 941-945. 

Turner, P.\,J., A.\,D.\,Thalerb, A.\,Freitag, and P.\,Colman Collins. 2019. Deep-sea hydrothermal vent ecosystem principles: Identification of ecosystem processes, services and communication of value. Marine Policy 101: 118-124.

Zhang, Y.\,S.\ and B.\ Silliman. 2019. A facilitation cascade enhances local biodiversity in seagrass beds. Diversity 11. doi:10.3390/d11030030


\end{hangparas}

\end{document}